% =========================================================
% Proof: Additivity of Suprema
% Source: volume-iii/analysis/bounding/notes/notes-bounds-arithmetic.tex
% =========================================================

\subsection*{Additivity of Suprema}
\label{prf:additivity-of-supremums}

\begin{remark*}[Return]
$\leftarrow$ \hyperref[thm:sup-add]{Back to Theorem in Notes}
\end{remark*}

\begin{theorem*}[Additivity of Suprema]
Let $A,B \subseteq \mathbb{R}$ be nonempty and bounded above. Define
\[
A+B=\{a+b : a\in A,\; b\in B\}.
\]
Then
\[
\sup(A+B)=\sup A+\sup B.
\]
\end{theorem*}

\begin{proof}
Let $s=\sup A$ and $t=\sup B$.

First we show that $s+t$ is an upper bound of $A+B$.

Let $a\in A$ and $b\in B$. Since $s=\sup A$ and $t=\sup B$, we have
\[
a \le s \quad \text{and} \quad b \le t.
\]
Adding these inequalities gives
\[
a+b \le s+t.
\]
Thus $s+t$ is an upper bound of $A+B$. Hence
\[
\sup(A+B) \le s+t.
\]

Next we show that $s+t$ is the least upper bound.

Let $\varepsilon>0$. Since $s=\sup A$, there exists $a\in A$ such that
\[
s-\frac{\varepsilon}{2} < a \le s.
\]
Similarly, since $t=\sup B$, there exists $b\in B$ such that
\[
t-\frac{\varepsilon}{2} < b \le t.
\]

Adding these inequalities gives
\[
s+t-\varepsilon < a+b \le s+t.
\]

Since $a+b \in A+B$, this shows that for every $\varepsilon>0$ there exists
$c\in A+B$ such that
\[
s+t-\varepsilon < c.
\]

Therefore, by the $\varepsilon$-characterization of the supremum,
\[
\sup(A+B)=s+t.
\]
\end{proof}

\begin{remark*}[Proof shape]
Two-part argument: exhibit the candidate value as an upper (lower) bound,
then show no smaller (larger) value qualifies using the
$\varepsilon$-characterization of the supremum (infimum).
\end{remark*}

\begin{remark*}[Dependencies]
Definition of supremum (infimum); \hyperref[prop:eps-char]{$\varepsilon$-Characterization, Proposition~\ref*{prop:eps-char}}.
\end{remark*}
